\section{Planteamiento del problema}
\label{cap:problemStatement}

\subsection{Tecnologías y Brechas del Estado del Arte}
\label{sec:literature_gap}

El análisis computacional de la alineación entre programas académicos y demandas del mercado laboral ha evolucionado significativamente en su representación textual. Los enfoques iniciales utilizaron modelos léxicos tradicionales como TF-IDF para ponderar palabras clave en descripciones de cursos; sin embargo, estas representaciones no logran capturar relaciones semánticas complejas ni sinonimias al depender estrictamente de frecuencias léxicas \citep{naseem_comprehensive_2021}. Para superar esto, se adoptaron representaciones distribuidas del lenguaje (\textit{word embeddings} como Word2Vec, GloVe y FastText), que ubican las palabras en espacios vectoriales donde la cercanía denota afinidad conceptual \citep{naseem_comprehensive_2021}. La evidencia empírica muestra que los \textit{word embeddings} superan a la bolsa de palabras, logrando medidas F de alineación mayores (0.571 frente a 0.333) \citep{miura_automated_2021}, y que clasificadores supervisados como SVM alcanzan exactitudes sobresalientes (85.35\%) en la extracción de capacidades \citep{fan_identifying_2023}.

En la actualidad, los modelos de lenguaje basados en transformadores multilingües (\textit{Sentence-Transformers} como SBERT/MPNet) permiten proyectar de manera conjunta textos en español e inglés dentro de un mismo espacio vectorial latente continuo $\mathbb{R}^d$, resolviendo las barreras idiomáticas y capturando la semántica contextual profunda \citep{cristea_course_2025}. Para medir la alineación en este espacio vectorial compartido, la similitud del coseno ($\cos \theta$) y las tasas de cobertura se han consolidado como métricas cuantitativas estándar \citep{miura_automated_2021}. Asimismo, el uso coordinado de algoritmos de agrupamiento no supervisado (\textit{clustering} K-Means, agrupamiento jerárquico y DBSCAN) permite descubrir la estructura disciplinar implícita y aislar asignaturas atípicas sin requerir etiquetado previo, evaluando la calidad geométrica mediante las métricas de Silhouette y Davies-Bouldin \citep{remadi_prompt_2024,fan_identifying_2023}.

Respecto a los antecedentes de la literatura comparada (Anexo B), destaca el trabajo de Russell et al. \citep{russell_e_mapping_2024}, quienes desarrollaron un mapeo computacional de planes de estudio hacia ocupaciones de la industria. No obstante, la solución propuesta en este escrito presenta diferencias sustanciales frente al trabajo de Russell et al.:
\begin{enumerate}[label=\alph*)]
	\item \textbf{Cobertura y Especificidad del Corpus:} Mientras Russell et al. limitaron su análisis a 2 instituciones de EE.UU. en el área acotada de Ciencia de Datos, este trabajo aborda 16 programas universitarios de educación superior en Chile (747 asignaturas), considerando el contexto multilingüe en español y la dualidad de mallas de ciclo largo (Civil) y ejecutoras.
	\item \textbf{Profundidad de Representación Contextual:} Russell et al. emplearon reglas léxicas y ponderación TF-IDF sobre descripciones breves de cursos (SkillNER). En contraste, nuestro método aplica transformadores vectoriales contextuales (SBERT) sobre mallas curriculares y perfiles de egreso declarados.
	\item \textbf{Triangulación Tripartita Simultánea:} A diferencia del mapeo unidireccional de Russell et al. hacia la base O*NET, nuestro método realiza una triangulación simultánea en el mismo espacio vectorial: \textit{Mallas Curriculares} $\leftrightarrow$ \textit{Estándares ACM/IEEE} $\leftrightarrow$ \textit{Ofertas de LinkedIn}.
\end{enumerate}

\subsection{Brecha Tecnológica}
\label{sec:technology_gap}

En el ecosistema universitario chileno no existe una plataforma automatizada implantada que evalúe y compare objetivamente los programas de estudio. Los rediseños curriculares a nivel local dependen de extensas reuniones de comisiones académicas manuales \citep{contreras_diseno_nodate} o de evaluaciones puntuales por pares evaluadores externos en procesos de acreditación \citep{diaz-quezada_rediseno_2019}. Por su parte, la medición de brechas de pertinencia laboral se limita a encuestas de percepción cualitativas aplicadas a egresados y empleadores, propensas al sesgo subjetivo humano \citep{latuz_pina_2024}. A nivel metodológico, las soluciones existentes en la literatura suelen corresponder a flujos experimentales diseñados ad hoc, basados en diccionarios estáticos sin profundidad semántica \citep{constantinov_preliminary_2015} o en modelos que demandan un costoso etiquetado manual de datos \citep{fan_identifying_2023}. Ninguna herramienta integra en una misma arquitectura en español la ingesta de mallas curriculares, el mapeo continuo contra la ACM y el contraste estratificado con el mercado laboral en un dashboard interactivo.

\subsection{Enunciado y Criterios del Problema}
\label{sec:problem_statement}

A partir de la brecha identificada, se formula la siguiente \textbf{pregunta de investigación}:
\begin{quote}
	\textit{“¿En qué medida difieren cuantitativamente la distribución disciplinar y la estructura de las mallas curriculares de Ingeniería Informática en Chile, y cuál es el grado de alineación semántica de sus asignaturas y perfiles de egreso declarados respecto a las directrices de los estándares internacionales ACM/IEEE (CS2013 y CC2020) y a las competencias tecnológicas demandadas por el mercado laboral local extraídas de LinkedIn?”}
\end{quote}

Se concibe como \textbf{conjetura o hipótesis de trabajo} que las mallas curriculares de corte tradicional de ciclo largo (Ingeniería Civil Informática) presentan una alta alineación semántica con los contenidos fundamentales del estándar ACM/IEEE pero manifiestan brechas de cobertura frente a competencias aplicadas de alta demanda en LinkedIn, mientras que las mallas curriculares de corte tecnológico o de ejecución muestran una mayor cercanía vectorial con las vacantes del mercado laboral local pero exhiben menor densidad en las bases teóricas del estándar internacional.

Para abordar esta problemática con rigor técnico y delimitación temporal, el estudio establece los siguientes criterios metodológicos y operacionales:
\begin{enumerate}[label=\arabic*.]
	\item \textbf{Delimitación de las Mallas del Corpus:} Se incluyen programas de formación universitaria vigentes y dictados oficialmente en Chile en el período 2018–2026, considerando tanto mallas de ciclo largo de 5.5 a 6 años como mallas de ejecución de 4 a 4.5 años (Anexo C).
	\item \textbf{Definición de "Diferencias Estructurales de Mallas":} El análisis de variaciones estructurales se evalúa cuantitativamente a través de dos ejes principales: (a) el \textit{Eje de Cobertura Disciplinar} (proporción de créditos y asignaturas de la malla asociados a las distintas \textit{Knowledge Areas} de la ACM), y (b) el \textit{Eje de Balance Modular} (relación entre el núcleo disciplinar obligatorio, líneas de especialización y materias de formación general).
	\item \textbf{Mecanismo de Relación y Estratificación con LinkedIn:} La demanda laboral se modela vectorizando las secciones de requisitos y funciones de 200 vacantes publicadas en LinkedIn durante 2024–2026. Para analizar el mercado de forma representativa, la muestra se divide formalmente en dos estratos: empresas del \textit{Sector Tecnológico (Mundo TI)} y empresas del \textit{Sector No-TI} (industria usuaria: banca, retail, minería, salud). La alineación se calcula evaluando la similitud del coseno $\cos(\mathbf{v}_{\text{asig}}, \mathbf{v}_{\text{job}})$ entre los vectores de las asignaturas de la malla y los requerimientos del mercado.
	\item \textbf{Criterios Numéricos de Decisión:} Se establecen como umbrales de validez: i) \textit{Alineación Normativa Aceptable}: Similitud de coseno promedio $\ge 0.70$ con las \textit{Knowledge Areas} de la ACM y Pureza de agrupamiento $\ge 0.70$; y ii) \textit{Presencia de Brecha Laboral Crítica}: Distancia semántica respecto a las competencias de LinkedIn $> 0.40$ (equivalente a similitud $\cos \theta < 0.60$).
\end{enumerate}